% ID: SEL-D0045-OA-E04
% ESTADO: cotejada. Segunda pasada visual y borrador compilado revisados.
% Fuente: Selectividad/sel_2014_matematicas_aplicadas/EXAMEN -ANDAL-5.pdf
% SHA-256: fe436fc3af8e56d0694316005aedeca2fffefbeb445aa9f60efa0aa6da4cba25
% Página: 1; opción A; ejercicio original 4.

Los responsables de tráfico de una ciudad trabajan con la hipótesis de que, al menos, el 65\% de sus habitantes son favorables a la creación de una red de carril-bici en esa ciudad.

Encuestados 950 habitantes, elegidos al azar, 590 están a favor de tal medida.
\begin{enumerate}[label=\alph*)]
\item \textbf{(1.5 puntos)} Mediante un contraste de hipótesis, $(H_0:p\geq0.65)$, con un nivel de significación del 10\%, \textquestiondown{}se puede decir que tienen razón los responsables de tráfico de esa ciudad?
\item \textbf{(1 punto)} \textquestiondown{}Se concluiría lo mismo si el nivel de significación fuera del 1\%?
\end{enumerate}
